\documentclass[12pt,a4paper]{article}

% --------------------
% Packages
% --------------------
\usepackage[margin=1in]{geometry}
\usepackage{graphicx}
\usepackage{fancyhdr}
\usepackage{setspace}
\usepackage{titlesec}
\usepackage{enumitem}
\usepackage{xcolor}
\usepackage{hyperref}
\usepackage{subcaption}
\usepackage[table]{xcolor}
\usepackage{colortbl}
\usepackage{makecell}
\usepackage{svg}

% --------------------
% Paragraph formatting
% --------------------
\setlength{\parindent}{0pt}  % Remove paragraph indentation
% \setlength{\parskip}{1em}    % Optional: add space between paragraphs

% --------------------
% Project Information
% --------------------
\newcommand{\CompanyName}{Pressure Control Solutions}
\newcommand{\ReportTitle}{Quattroflow Pressure Dynamics}
\newcommand{\ReportSubTitle}{Without and With an Downstream Equilibar}
\newcommand{\AuthorName}{Hubald Verzijl}
\newcommand{\AuthorEmail}{hubald@pressurecontrolsolutions.com}
\newcommand{\VersionNumber}{1.1}
\newcommand{\ReportDate}{\today}

\hypersetup{
    colorlinks=true,
    linkcolor=blue,
    filecolor=magenta,      
    urlcolor=cyan,
    pdftitle={\ReportTitle},
    pdfauthor={\AuthorName},
    pdfkeywords={Pressure, Dynamics, Quattroflow, Equilibar, Report},
}

% --------------------
% Header & Footer
% --------------------
\pagestyle{fancy}
\fancyhf{} % Clear default header/footer

% Header
\setlength{\headheight}{30pt}
\fancyhead[L]{%
  \includegraphics[height=0.75cm]{images/logo.png}
}
\fancyhead[R]{Report Version \VersionNumber}

% Footer
\fancyfoot[C]{\thepage}

\renewcommand{\headrulewidth}{0.4pt}
\renewcommand{\footrulewidth}{0.4pt}

% --------------------
% Title Page
% --------------------
\begin{document}

\begin{titlepage}
    \thispagestyle{fancy}
    \fancyfoot{} % Clear footer for this page only
    \centering
    \vspace*{2cm}
    {\Huge \bfseries \ReportTitle \par}
    \vspace{0.5cm}
    {\Large \ReportSubTitle \par}
    \vspace{2cm}
    \begin{center}
        \includegraphics[height=7.5cm]{images/qf1200su-ht.jpg}
    \end{center}
    \vspace{2cm}
    {\Large \CompanyName \par}
    \vspace{0.25cm}
    {\small Author: \AuthorName \par}
    \vspace{0.25cm}
    {\small Report Version: \VersionNumber \par}
    \vspace{0.25cm}
    {\small Date: \ReportDate \par}
    \vfill
\end{titlepage}

% --------------------
% Table of Contents
% --------------------
\section*{Introduction}
This short report describes the pressure dynamics of a Quattroflow \textcolor{red}{xxx} pump, both without and with an Equilibar backpressure regulator installed downstream of the pump. The Quattroflow pump is operated at different speeds ranging from 250 to 1000 RPM. As the roller pump contains 4 diaphragm pumps that are sequential driven by an excentric shaft, 4 pressure pulses are expected per revolution of the pump. As the experiments are performed without and with an Equilibar backpressure regulator installed downstream of the pump, it is expected that the Equilibar will significantly reduce the pressure fluctuations upstream of the regulator generated by the pump.

\tableofcontents
\newpage

\section{Method}
To model a customer proces, a needle valve (Fitok MH6S-FNS4-V) is installed downstream of the proces to create a certain backpressure. For the first part, no Equilibar is installed, and both pressure sensors are connected directly to the outlet of the pump (in serie or parallel). For the second part, an Equilibar is installed downstream of the pump, and a different backpressures (0.25 to 2.5 bar) and pump rotations (250, 500, 750 and 1000 RPM) are applied. Pressure sensors are placed up- and downstream of the Equilibar. Figures \ref{fig:setup_with_equilibar} shows the setup with Equilibar. \\
During the experiments, two Keller 20SX sensors are used and connected to a PCS EQ-DA PCB for data acquisition. 
Each sensor is sampled at 1250 Hz and filtered with a digital biquad filter with a cutoff frequency of 1000 Hz. All collected data, including MCU timestamps, is communicated to a custom-designed GUI for visualization and logging to a \texttt{.csv} file. Further analysis is performed using \texttt{Python}, where a small portion of the signals (0.2 seconds) is shown on the x-axis throughout this report. As different absolut values might be large compared to the relative fluctuations, in some figures the average of the signals is subtracted. These pressures are called \textit{pressure fluctuation} $p'(t)$ and are determined by
\begin{equation}
    p'(t) = p(t) - \bar{p},
\end{equation}
where $p(t)$ is the absolute pressure signal and $\bar{p}$ is the average pressure over the shown time interval.

\begin{figure}[h!]
    \centering
    \includegraphics[width=0.8\textwidth]{images/setup.jpeg}
    \caption{Setup with Equilibar. Quattroflow pump is shown top left, Equilibar downstream of pump with up- and downstream pressure sensors connected to a EQ-DA for data acquisition, and needle valve at the end of the setup to create backpressure. The reference pressure at the Equilibar is controlled by a second EQ-DA shown at the bottom.}
    \label{fig:setup_with_equilibar}
\end{figure}




% --------------------
% Report Content
% --------------------
\section{Results without Equilibar}
\subsection{Sensor validation}
Before performing the actual experiments, both Keller 20SX sensors are validated by connecting them in serie and in parallel downstream of the Quattroflow pump, which is shown in Figure \ref{fig:serie} and \ref{fig:parallel}. The needle valves resembles the userprocess resistance. 

\begin{figure}[h!]
    \centering
    \begin{minipage}{0.45\textwidth}
        \centering
        \includegraphics[width=\textwidth]{images/serie.jpeg}
        \caption{Sensors connected in serie for validation test.}
        \label{fig:serie}
    \end{minipage}
    \hfill
    \begin{minipage}{0.45\textwidth}
        \centering
        \includegraphics[width=\textwidth]{images/parallel.jpeg}
        \caption{Sensors connected in parallel for validation test.}
        \label{fig:parallel}
    \end{minipage}
    \hfill
\end{figure}

The results of this validation test are shown in Figure \ref{fig:sensor_validation}. From Figure \ref{fig:validation_serie}, it is clear that an increase in RPMs while using a needle valve downstream leads to an increased absolute pressure, as expected. The relative fluctuations are shown in Figure \ref{fig:validation_serie_mean_subtracted}. However, the up- and downstream pressures in the series configuration are not equal, while this is the case in the parallel configuration (Figure \ref{fig:validation_parallel_mean_subtracted}). Although this indicates that the sensors are measuring correctly, it also shows that the fluidic system introduces some oscillatory behavior, especially at higher frequencies, leading to some phase differences and amplitude differences between both sensors in the series configuration.

\begin{figure}[htbp!]
    \begin{subfigure}{\linewidth}
        \centering
        \includegraphics[width=\textwidth]{images/2_serie_open.eps}
        \caption{Sensor validation test with sensors connected in \textbf{serie} (Figure \ref{fig:serie}).}
        \label{fig:validation_serie}
    \end{subfigure}

    \par\bigskip

    \begin{subfigure}{\linewidth}
        \centering
        \includegraphics[width=\textwidth]{images/2_serie_open_mean_subtracted.eps}
        \caption{Sensor validation test with sensors connected in \textbf{serie} (Figure \ref{fig:serie}). Pressure fluctuations with average subtracted.}
        \label{fig:validation_serie_mean_subtracted}
    \end{subfigure}

    \par\bigskip

    \begin{subfigure}{\linewidth}
        \centering
        \includegraphics[width=\textwidth]{images/1_parallel_open_mean_subtracted.eps}
        \caption{Sensor validation test with sensors connected in \textbf{parallel} (Figure \ref{fig:parallel}). Pressure fluctuations with average subtracted.}
        \label{fig:validation_parallel_mean_subtracted}
    \end{subfigure}

    \caption{Results of the sensor validation tests. Both sensors show similar pressure dynamics in the \textbf{parallel} case, indicating that they are functioning correctly.}
    \label{fig:sensor_validation}
\end{figure}

\subsection{Proces pressure}
The needle valve creates a resistance in the fluidic system, leading to a certain backpressure/proces pressure depending on the pump speed. The absolute upstream and downstream pressures at different pump speeds are shown in Figure \ref{fig:validation_serie}. From this figure, it can be seen that the pressure increases with increasing pump speed, as expected. Table \ref{tab:proces_pressures} summarizes the average proces pressures at different pump speeds.

\begin{table}[h!]
    \centering
    \begin{tabular}{|c|c|c|c|}
        \hline
        \textbf{\makecell{Speed \\ {[RPM]}}} 
        & \textbf{\makecell{Proces pressure \\ {[mbar]}}} \\
        \hline
        250  & 0 \\
        \hline
        500  & 250 \\
        \hline
        750  & 500 \\
        \hline
        1000 & 750 \\
        \hline
    \end{tabular}
    \caption{Average proces pressures at different pump speeds created by the needle valve.}
    \label{tab:proces_pressures}
\end{table}


\newpage
\section{Results with Equilibar}
The setup from Figure \ref{fig:setup_with_equilibar} is used to perform experiments with an Equilibar downstream of the Quattroflow pump. Pressure measurement are performed at variable backpressures (0, 250, 500, 750, 1000, 1250, 1500, 1750, 2000, 2250, and 2500 mbar) and at different pump speeds (250, 500, 750, and 1000 RPM). \\
Important to note is that when the RPMs increase, a process pressure (Table \ref{tab:proces_pressures}) can be genreated that is higher than the impoased reference pressre which leads to a fully opeened diaphragm (Figure \ref{fig:equilibar_variable_backpressure_results_low}). When the reference pressure is sufficiently above the proces pressure, the diaphragm starts to close, leading to periodic opening and closing behavior (Figure \ref{fig:equilibar_variable_backpressure_results_high}).

\begin{figure}[h!]
    \centering
    \includegraphics[width=1\textwidth]{images/6_eq_500mbar_open.eps}
    \caption{Pressure dynamics with an Equilibar at variable pump speeds. At higher pump speeds, the process pressure exceeds the imposed reference pressure. In this case, the diaphragm is fully open.}
    \label{fig:equilibar_variable_backpressure_results_low}
\end{figure}

\begin{figure}[h!]
    \centering
    \includegraphics[width=1\textwidth]{images/6_eq_1500mbar_open.eps}
    \caption{Pressure dynamics with an Equilibar at variable pump speeds. At higher pump speeds, the imposed reference pressure is still higher dan the process pressure. In all situations, this diaphragm is periodically open.}
    \label{fig:equilibar_variable_backpressure_results_high}
\end{figure}


\newpage

\section{Result interpretation}
\subsection{Maximum backpressure (2500 mbar)}
To quantify the effect of the Equilibar on pressure fluctuations upstream and downstream, an analysis is performed to evaluate the percentage change in pressure fluctuations relative to the configuration without an Equilibar. For both configurations, the root mean square (RMS) value of the upstream and downstream pressure signals is calculated and defined as
\begin{equation}
p_{\mathrm{RMS}} = \sqrt{\frac{1}{N} \sum_{i=1}^{N} \left(p'(t_i)\right)^2},
\end{equation}
where $p'(t_i)$ denotes the pressure fluctuation at time $t_i$ and $N$ is the total number of samples.
Subsequently, the percentage change in RMS pressure fluctuations is determined. In this analysis, the pressure fluctuations measured at a backpressure of 2500 mbar with the Equilibar are compared to those obtained without an Equilibar. The results are summarized in Table~\ref{tab:pressure_changes}. The data show that the use of the Equilibar leads to a \textbf{reduction in upstream pressure fluctuations}, while \textbf{downstream pressure fluctuations increase}.

\begin{table}[h!]
    \centering
    \begin{tabular}{|c|c|c|c|}
        \hline
        \textbf{\makecell{Speed \\ {[RPM]}}} 
        & \textbf{\makecell{Backpressure \\ {[mbar]}}}
        & \textbf{\makecell{Upstream change \\ {[\% RMS]}}}
        & \textbf{\makecell{Downstream change \\ {[\% RMS]}}} \\
        \hline
        250  & 2500 & \cellcolor{green!22}-22\% & \cellcolor{red!44}44\% \\
        \hline
        500  & 2500 & \cellcolor{green!48}-48\% & \cellcolor{red!27}27\% \\
        \hline
        750  & 2500 & \cellcolor{green!48}-48\% & \cellcolor{red!16}16\% \\
        \hline
        1000 & 2500 & \cellcolor{green!43}-43\% & \cellcolor{red!19}19\% \\
        \hline
    \end{tabular}
    \caption{Percentage change in pressure fluctuations [\% RMS] for upstream and downstream at constant backpressure of 2500 mbar. Comparison between without and with Equilibar cases. Green indicates a decrease in fluctuations, red indicates an increase in fluctuations.}
    \label{tab:rms_pressure_changes}
\end{table}

\subsection{Optimal backpressure}
From previous experiments, it was found that the optimal backpressures for minimizing upstream and downstream pressure fluctuations were approximately 0, 250, 500, and 750 mbar for pump speeds of 250, 500, 750, and 1000 RPM, respectively. A similar analysis as in the previous section is performed to evaluate the percentage change in pressure fluctuations relative to the configuration without an Equilibar. The results are summarized in Table~\ref{tab:pressure_changes_optimal}. The data show that the use of the Equilibar at optimal backpressures leads to a \textbf{significant reduction in both upstream and downstream pressure fluctuations}.

\begin{table}[h!]
    \centering
    \begin{tabular}{|c|c|c|c|}
        \hline
        \textbf{\makecell{Speed \\ {[RPM]}}} 
        & \textbf{\makecell{Backpressure \\ {[mbar]}}}
        & \textbf{\makecell{Upstream change \\ {[\% RMS]}}}
        & \textbf{\makecell{Downstream change \\ {[\% RMS]}}} \\
        \hline
        250  & 0 & \cellcolor{green!71}-71\% & \cellcolor{green!69}-69\% \\
        \hline
        500  & 250 & \cellcolor{green!65}-65\% & \cellcolor{green!48}-48\% \\
        \hline
        750  & 500 & \cellcolor{green!38}-38\% & \cellcolor{green!67}-67\% \\
        \hline
        1000 & 750 & \cellcolor{green!35}-35\% & \cellcolor{green!47}-47\% \\
        \hline
    \end{tabular}
    \caption{Percentage change in pressure fluctuations [\% RMS] for upstream and downstream at optimal backpressures. Comparison between without and with Equilibar cases. Green indicates a decrease in fluctuations, red indicates an increase in fluctuations.}
    \label{tab:rms_pressure_changes_optimal}
\end{table}

\clearpage
\section{Equilibar in parallel}
A setup is created where are Equilibar is set in parallel. The reference pressure on the Equilibar is created by an EQ-DA. A ball valve is used to model the process. During the experiments, the valve is controled manually, such that the upstream pressure is around 1 bar. This setup has potentially two advantage:
\begin{itemize}
    \item \textbf{Limiting maximum process pressure:} If the resistance in the process increases, the membrane pump will not lower the RPMs automatillcy, so the upstream pressure can increase significantly. By placing the Equilibar in parallel, the maximum process pressure can be limited to the applied reference pressure at the Equilibar, which can potentially protect the system from overpressure.
    \item \textbf{Damping pressure fluctuations:} The Equilibar can potentially dampen the pressure fluctuations generated by the pump, leading to a more stable pressure output. This can be especially beneficial in applications where precise pressure control is required.
\end{itemize}


\begin{figure}[h!]
  \centering
  \includesvg[width=1\textwidth]{images/20260206/Setup_1.svg}
  \label{fig:setup_parallel}
  \caption{Setup with Equilibar in parallel. Reservoir shown in red, Quattroflow pump shown in green, Equilibar with EQ-DA shown in orange, and pressure sensors shown in black. The process is modeled by a ball valve shown in blue. In the dotted pliers can be placed to disable the flow through the Equilibar}
\end{figure}

To prove the second advantage, some experiments are perfomred at 250, 500, 750 and 1000RPM. The results are shown in Figure \ref{fig:parallel_results} and \ref{fig:parallel_rms_vs_rpm}. From these figures, it can be seen that the pressure fluctuations are significantly reduced when using the Equilibar in parallel, especially at lower pump speeds. This indicates that the Equilibar is effective in damping the pressure fluctuations generated by the Quattroflow pump. 

% INlcude figure
\begin{figure}[h!]
    \centering
    \includegraphics[width=0.8\textwidth]{images/20260206/pressure_vs_time_1_5ref.eps}
    \caption{Pressure dynamics with an Equilibar in parallel at variable backpressures.}
    \label{fig:parallel_results}
\end{figure}

\begin{figure}
    \centering
    \includegraphics[width=0.8\textwidth]{images/20260206/rms_pressure_vs_rpm.eps}
    \caption{Pressure dynamics with an Equilibar in parallel at variable backpressures. Pressure fluctuations with average subtracted.}
    \label{fig:parallel_rms_vs_rpm}
\end{figure}

\clearpage
\newpage


\begin{table}[h!]
    \centering
    \begin{tabular}{|
        c
        |c
        |>{\centering\arraybackslash}m{2cm}
        |>{\centering\arraybackslash}m{2cm}
        |>{\centering\arraybackslash}m{2cm}
        |>{\centering\arraybackslash}m{2cm}|
    }
        \hline
        \textbf{RPM} 
        & \textbf{Backpressure}
        & \multicolumn{2}{c|}{\textbf{Upstream change}} 
        & \multicolumn{2}{c|}{\textbf{Downstream change}} \\
        \cline{3-6}
        &
        & \textbf{Base} & \textbf{4x} 
        & \textbf{Base} & \textbf{4x} \\
        \hline
        250  
        & 2500 mbar
        & \cellcolor{green!20}-17\%
        & \cellcolor{green!40}-39\%
        & \cellcolor{red!70}73\%
        & \cellcolor{green!50}-51\% \\
        \hline
        500  
        & 2500 mbar
        & \cellcolor{green!50}-52\%
        & \cellcolor{green!30}-27\%
        & \cellcolor{red!50}46\%
        & \cellcolor{green!40}-42\% \\
        \hline
        750  
        & 2500 mbar
        & \cellcolor{green!60}-55\%
        & \cellcolor{red!100}188\%
        & \cellcolor{red!30}29\%
        & \cellcolor{red!10}5\% \\
        \hline
        1000
        & 2500 mbar
        & \cellcolor{green!40}-44\%
        & \cellcolor{green!50}-54\%
        & \cellcolor{red!50}47\%
        & \cellcolor{green!80}-80\% \\
        \hline
    \end{tabular}
    \caption{Percentage change in pressure fluctuations for upstream and downstream at base and 4x frequencies. Comparison between without and with Equilibar (at 2500 mbar) cases. Green indicates a decrease in fluctuations, red indicates an increase in fluctuations.}
    \label{tab:pressure_changes}
\end{table}

\begin{table}[h!]
    \centering
    \begin{tabular}{|
        c
        |c
        |>{\centering\arraybackslash}m{2cm}
        |>{\centering\arraybackslash}m{2cm}
        |>{\centering\arraybackslash}m{2cm}
        |>{\centering\arraybackslash}m{2cm}|
    }
        \hline
        \textbf{RPM} 
        & \textbf{Backpressure}
        & \multicolumn{2}{c|}{\textbf{Upstream change}} 
        & \multicolumn{2}{c|}{\textbf{Downstream change}} \\
        \cline{3-6}
        & 
        & \textbf{Base} & \textbf{4x} 
        & \textbf{Base} & \textbf{4x} \\
        \hline
        250
        & 0 mbar
        & \cellcolor{green!80}-77\%
        & \cellcolor{green!80}-79\%
        & \cellcolor{green!70}-74\%
        & \cellcolor{green!80}-76\% \\
        \hline
        500
        & 250 mbar
        & \cellcolor{green!80}-80\%
        & \cellcolor{red!30}27\%
        & \cellcolor{green!40}-37\%
        & \cellcolor{green!60}-57\% \\
        \hline
        750
        & 500 mbar
        & \cellcolor{green!90}-88\%
        & \cellcolor{red!100}583\%
        & \cellcolor{green!80}-74\%
        & \cellcolor{green!10}-7\% \\
        \hline
        1000
        & 750 mbar
        & \cellcolor{green!80}-75\%
        & \cellcolor{red!40}40\%
        & \cellcolor{green!50}-51\%
        & \cellcolor{green!60}-61\% \\
        \hline
    \end{tabular}
    \caption{Percentage change in pressure fluctuations for upstream and downstream at base and 4x frequencies. Comparison between without and with Equilibar (at optimal backpressures) cases. Green indicates a decrease in fluctuations, red indicates an increase in fluctuations.}
    \label{tab:pressure_changes_optimal}
\end{table}


    
\newpage


\section{Results without Equilibar}
For the first experiment, both sensors are connected directly to the outflow of the Quattroflow pump as shown in Figure \ref{fig:setup_1}. 
\begin{figure}[h!]
    \centering
    \includegraphics[width=0.5\textwidth]{images/no_equilibar.png}
    \caption{An EQ-DA is used to read out the two Keller 20SX sensors.}
    \label{fig:setup_1}
\end{figure}

Several tests are performed at different pump speeds: 250, 500, 750, and 1000 RPM. This results in approximately the following expected pulsing frequencies: 17, 33, 50, and 67 Hz.

\begin{figure}[h!]
    \centering
    \includegraphics[width=1\textwidth]{images/without_equilibar.eps}
    \caption{Pressure dynamics without an Equilibar.}
    \label{fig:without_equilibar_results}
\end{figure}

From Figure \ref{fig:without_equilibar_results}, it can be seen that at lower RPMs (250 and 500 RPM) both pressure curves look similar, except for some absolute pressure differences. This can be due to the fact that the sensors were not zeroed before the experiment. Additionally, a phase shift can be observed, which can be explained by the physical distance between both sensors. At higher RPMs (750 and 1000 RPM), the absolute pressure increases slightly, and the amplitude of the pressure fluctuations increases on the upstream side. The downstream side shows smaller pressure fluctuations at these higher frequencies. This can be explained by the damping effect of the differential pressure across the first sensor or some low-pass filtering within the sensor itself. It would be interesting to interchange the sensors to determine whether the effect is sensor-related or location-related; unfortunately, this was not performed during this experiment.

\section{Results with Equilibar and constant backpressure}
An Equilibar is connected to the outlet of the pump. The first pressure sensor is installed directly at the outlet of the pump (upstream side of the Equilibar) and the second pressure sensor is installed downstream of the Equilibar (downstream side). An EQ-DA is used to control the reference pressure at the Equilibar. The design of the Equilibar results in an upstream pressure that approximately equals the applied reference pressure. During these experiments, a constant reference pressure of 1 bar is applied at the Equilibar. Subsequently, the pump is again operated at different speeds: 250, 500, 750, and 1000 RPM.

\begin{figure}[h!]
    \centering
    \includegraphics[width=1\textwidth]{images/with_equilibar.eps}
    \caption{Pressure dynamics with an Equilibar.}
    \label{fig:with_equilibar_results}
\end{figure}

In Figure \ref{fig:with_equilibar_results}, the pressure dynamics with an Equilibar are shown. The top graphs show the absolute pressure over time for both sensors at different pump speeds. The bottom graphs show an average-subtracted pressure signal to better visualize the (relative) pressure fluctuations. \\
From this figure, it can be seen that the upstream pressure (sensor 1) again shows the expected pressure fluctuations, which increase in amplitude with increasing pump speed. However, the downstream pressure (sensor 2) now shows a more stable pressure with minimal fluctuations. This indicates that the Equilibar is effective in damping the downstream pressure fluctuations generated by the Quattroflow pump, even at higher pump speeds (peak-to-peak 10 mbar vs 30 mbar). However, the upstream pressure fluctuations are not damped and are slightly larger compared to the results without an Equilibar (peak-to-peak 80 mbar vs 60 mbar). In Section \ref{sec:conclusion}, some possible explanations are given for this observation. \\
It is important to note that some damping was observed in the `without Equilibar' experiment as well, especially at higher pump speeds. This could be related to the sensor response time, filtering within the sensor or the damping effect of the fluidic system itself. Observed damping might therefore not only be an effect of the Equilibar but also of the fluidic system and/or sensors used.

The frequency content of the pressure signals provides more insight into the pressure dynamics. In Figure \ref{fig:with_equilibar_freq_results}, this is shown for both sensors at different pump speeds. The expected pulsing frequencies (17, 33, 50, and 67 Hz) are again clearly visible in both sensors. This analysis can be used to determine the damping effect.

\begin{figure}[h!]
    \centering
    \includegraphics[width=1\textwidth]{images/with_equilibar_freq.eps}
    \caption{Frequency representation of pressure dynamics with an Equilibar.}
    \label{fig:with_equilibar_freq_results}
\end{figure}

\subsection{Comparing the without and with Equilibar cases}
To compare the results of the `without Equilibar' and `with Equilibar' cases, the frequency representation of both experiments at 1000 RPM is shown in Figure \ref{fig:with_without_comparison_results}. It can be seen that the upstream pressure fluctuations are slightly increased when using the Equilibar compared to the setup without the Equilibar. A possible explanation for this observation is that the Equilibar introduces some additional flow resistance in the system, which could lead to increased pressure fluctuations upstream. This could be explained if the dynamics of the Equilibar are too slow to respond to the fast pressure pulses generated by the Quattroflow pump. Section \ref{sec:conclusion} elaborates a little bit further on this.\\
\begin{figure}[h!]
    \centering
    \includegraphics[width=1\textwidth]{images/with_without_comparison.eps}
    \caption{Comparing the pressure dynamics with and without an Equilibar at 1000 RPM.}
    \label{fig:with_without_comparison_results}
\end{figure}

\newpage
\section{Results with Equilibar and variable backpressures}
In this experiment, the Equilibar is again connected to the outlet of the pump with the same sensor configuration as in the previous section. However, the reference pressure at the Equilibar is varied between 0.25, 0.5, and 1 bar. The pump is operated constantly at a speed of 1000 RPM.

\begin{figure}[h!]
    \centering
    \includegraphics[width=1\textwidth]{images/with_equilibar_variable.eps}
    \caption{Pressure dynamics with an Equilibar at variable backpressures.}
    \label{fig:with_equilibar_variable_results}
\end{figure}

During analysis, the absolute pressure is plotted as well as a relative pressure signal with the average subtracted. From Figure \ref{fig:with_equilibar_variable_results}, it can be seen that the upstream pressure fluctuations increase slightly with increasing backpressure. This can be explained by the fact that at higher backpressure, the Quattroflow pump has to work against a higher load, resulting in slightly larger pressure fluctuations. However, the downstream pressure remains stable with smaller fluctuations, regardless of the applied backpressure. \\
Again, some of the damping effect could be related to the fluidic system and/or sensors used, as was also observed in the previous section.

\newpage
\section{Conclusion}
\label{sec:conclusion}
It was expected that the Equilibar would significantly reduce the pressure fluctuations upstream of the regulator that are generated by the pump. However, the results show that the upstream pressure fluctuations are slightly increased when using the Equilibar compared to the setup without the Equilibar. A possible explanation for this observation is that the Equilibar introduces some additional flow resistance in the system, which could lead to increased pressure fluctuations upstream. This could be explained if the dynamics of the Equilibar are too slow to respond to the fast pressure pulses generated by the Quattroflow pump. 

\bigskip

Based on the results and observations from this report, it can be concluded that the Equilibar is effective in damping the pressure fluctuations downstream of the pump, providing a more stable pressure output (and flow output). However, further investigation is needed to fully understand the impact of the Equilibar on upstream pressure dynamics.

\bigskip

Recommendations for further investigation:
\begin{itemize}
    \item Interchange the sensors to determine if the (damping) effect is sensor-related or location-related;
    \item Ensure no air bubbles are present in the fluidic system, as they could potentially influence the pressure readouts;
    \item Use a faster-responding Equilibar model to determine if the dynamics of the Equilibar are too slow to respond to the fast pressure pulses generated by the Quattroflow pump. Changing the diaphragm material could potentially help achieve this;
    \item Actively reduce fluctuations by closed-loop controlling the reference pressure at the Equilibar based on the fluctuating upstream pressure sensor signal. The downstream pressure could also potentially be controlled to maintain a even more constant pressure. The latter could also be used to create a (more) constant output flow.
\end{itemize}
% --------------------
% References
% --------------------
% \begin{thebibliography}{9}
% \bibitem{ref1}
% Author, \textit{Title of Reference}, Publisher, Year.
% \end{thebibliography}

% Punten
% Fluctuaties upstream of juist downstream verminderen?
% When process pressure is equal to reference pressure, the air in reference pressure acts as an compressible volume, which can dampen fluctuations?

% Vervolgstappen
% Overleg met Bas over huidige resultaten
% Wens: 
% - Flowcontrol zonder stabiele voordruk
% - Stabiele voordruk met continue afblazen? Of juist met actieve regeling?
% 13 februari naar PCS?

% 6 februari 
% Scan through all different resistances for variable frequencies




\end{document}
